%% ma: L11-B09-0016
%% lop: 11
%% chuong: 3
%% bai: 9
%% dang: 11.9.4
%% loai: TL
%% muc: TH
%% dap_an: "-"
%% nguon: "de-on-tap-giua-ki-1-toan-11-nam-2026-2027-truong-thpt-ngo-quyen-da-nang.pdf (M09, Đề ôn tập giữa kì 1 Toán 11 THPT Ngô Quyền, Đà Nẵng), Mã đề 0113, Phần IV Câu 3"
%% kien_thuc: "Bài 9 (các số đặc trưng đo xu thế trung tâm)"
%% trang_thai: chinh-thuc
\begin{ex}
Bảng dưới đây thống kê thành tích nhảy xa của một số học sinh lớp 12:
\begin{center}\begin{tabular}{|l|c|c|c|c|c|}\hline
Thành tích (cm) & $[150;180)$ & $[180;210)$ & $[210;240)$ & $[240;270)$ & $[270;300)$\\\hline
Số học sinh & $3$ & $5$ & $28$ & $14$ & $8$\\\hline
\end{tabular}\end{center}
Tính số cm trung bình nhảy xa của học sinh lớp 12 (kết quả làm tròn đến hàng phần mười).
\loigiai{$n=58$, giá trị đại diện $165;195;225;255;285$. \[\bar x=\dfrac{165\cdot3+195\cdot5+225\cdot28+255\cdot14+285\cdot8}{58}=\dfrac{13620}{58}\approx234{,}8\text{ (cm).}\]}
\end{ex}
